\section{幂级数是一族函数}
\label{sec:02-Calculus/05-Sequences-and-Series/05}

上一章我们反复用过这个看似无所不能的等式：
\[
\frac1{1-x}=1+x+x^2+x^3+\cdots.
\]

现在试着把 \(x=2\) 强行代入：左边是 \(\frac{1}{1-2}=-1\)；右边则是 \(1+2+4+8+\cdots = +\infty\)。一串正数累加，竟然等于 \(-1\)！

整个推导过程中的代数变换没有一步出错，错的是那个等号——\textbf{它从一开始就附带了一片``合法地盘'，只是边界没有写在等号上}。

把 \(x\) 限制在合法地盘内，等式就是真理；一旦跨出边界，代数形式就沦为毫无意义的废话。本节的任务，就是把这块地盘的边界精准地划出来，并解释为什么它一定是以展开中心为圆心的``对称安全区'。

\subsection{一个代数式，一族数项级数}

以 \(a\) 为中心的幂级数形如：
\[
\sum_{n=0}^{\infty}c_n(x-a)^n.
\]

严格来说，它本身还不是一个确定的级数，而是一个\textbf{``级数生成器'}。
只有当我们将 \(x\) 钉在数轴上的某个具体点时，它才塌缩成前几节讨论的那种数项级数。

换一个 \(x\)，就得到一个全新的级数，命运也截然不同。因此，一个幂级数打包携带了两个核心信息：
\begin{enumerate}
    \item \textbf{定义域}：哪些 \(x\) 能让生成的级数收敛？
    \item \textbf{映射规则}：在这些收敛的 \(x\) 上，极限值是多少？
\end{enumerate}
这正是一个函数的定义。只是它的``定义域'不能直接看出来，必须由我们亲自解出来。

\subsection{为什么收敛地盘一定是个对称区间？}

凭什么收敛区域不能左边宽、右边窄？凭什么它一定是围绕中心 \(a\) 的对称区间？

秘密藏在第 \(n\) 项的结构中：\(c_n(x-a)^n\) 的大小取决于两股力量的博弈——系数 \(c_n\) 的增长或衰减，以及偏离中心距离的累幂 \(|x-a|^n\)。

我们可以用\textbf{根值检验法}来测量它的收敛性。第 \(n\) 项开 \(n\) 次方后，我们考察其极限：
\[
L = \lim_{n\to\infty} \sqrt[n]{|c_n(x-a)^n|} = |x-a| \cdot \limsup_{n\to\infty}|c_n|^{1/n}.
\]
只要 \(L < 1\)，级数就绝对收敛；若 \(L > 1\)，项就会发散到无穷。

请注意：\textbf{这个判据对 \(x\) 的依赖，完全且仅通过距离 \(|x-a|\) 来体现！}
无论向左偏离还是向右偏离，只要距离满足 \(|x-a| < R\)，收敛性就完全一致。这就一脚锁死了一维空间里收敛地盘的几何形状——它只能是以 \(a\) 为中心、半径为 \(R\) 的对称区间 \((a-R, a+R)\)。

其中收敛半径 \(R\) 由著名的 Cauchy--Hadamard 公式给出：
\[
\frac1R=\limsup_{n\to\infty}\abs{c_n}^{1/n}.
\]
（在实际计算中，若系数比值极限存在，用比值检验算 \(\frac{1}{R} = \lim \abs{\frac{c_{n+1}}{c_n}}\) 通常更省事。）

当 \(|x-a| > R\) 时，级数彻底发散；而当 \(|x-a| = R\)（即两个边界端点 \(x = a \pm R\)）时，正好处于 \(L=1\) 的临界线，即上一节所说的``测量工具分辨率用尽'的位置。边界上的命运，必须单独逐一排查。

\begin{center}
\begin{tikzpicture}[x=2.6cm,y=1cm,font=\small,>=stealth]
  \fill[blue!12] (-1,-0.22) rectangle (1,0.22);
  \fill[red!8] (-1.75,-0.22) rectangle (-1,0.22);
  \fill[red!8] (1,-0.22) rectangle (1.75,0.22);
  \draw[->,black!60] (-1.9,0) -- (1.98,0);
  \node[below,black!60] at (1.93,-0.06) {\(x\)};
  \draw[black!70,thick] (-1,-0.34) -- (-1,0.34);
  \draw[black!70,thick] (1,-0.34) -- (1,0.34);
  \fill (0,0) circle (1.7pt);
  \node[below] at (0,-0.4) {\(0\)};
  \node[below] at (-1,-0.4) {\(-1\)};
  \node[below] at (1,-0.4) {\(1\)};
  \node[black!75] at (0,0.58) {\textbf{绝对收敛区} (\(|x| < 1\))};
  \node[black!60,font=\footnotesize] at (-1.42,0.58) {发散区};
  \node[black!60,font=\footnotesize] at (1.42,0.58) {发散区};
  \draw[->,black!55] (-1.0,1.3) -- (-1.0,0.42);
  \node[above,align=center,black!75,font=\footnotesize] at (-1.0,1.3) {\textbf{左端点收敛}\\（交错调和级数）};
  \draw[->,black!55] (1.0,1.3) -- (1.0,0.42);
  \node[above,align=center,black!75,font=\footnotesize] at (1.0,1.3) {\textbf{右端点发散}\\（正项调和级数）};
\end{tikzpicture}

\emph{以 \(\sum \frac{x^n}{n}\) 为例，收敛半径 \(R=1\)。内部由指数级衰减一刀切开，强行保障收敛；但到了边界 \(|x|=1\) 处，指数力量衰减完毕，端点的生死存亡完全取决于代数尺度的衰减速度以及正负抵消机制。}
\end{center}

同样是 \(R=1\)，不同级数在边界端点的表现千差万别：
\begin{itemize}
    \item \(\sum x^n\)：在 \(x = \pm 1\) 两端都发散；
    \item \(\sum \frac{x^n}{n}\)：在 \(x = -1\)（交错调和）收敛，在 \(x = 1\)（调和级数）发散；
    \item \(\sum \frac{x^n}{n^2}\)：在 \(x = \pm 1\) 两端都绝对收敛。
\end{itemize}

这里的深层逻辑在于：\textbf{收敛半径由系数的``指数尺度'决定，而端点生死由系数的``代数多项式尺度'决定}。指数尺度统治内部，代数尺度与符号抵消统治边界，这正是上一节``两把尺子'分工的再次体现。

\subsection{几何级数：所有幂级数的母体}

\[
\sum_{n=0}^{\infty}x^n=\frac1{1-x},\qquad\abs x<1
\]
几何级数不仅是最简单的幂级数，更是整个幂级数家族的``原版母体'。通过对它进行换元、求导和积分，我们可以像变形金刚一样衍生出许许多多复杂函数的展开式。

只要限定在收敛区间内部（\(|x| < 1\)），这些微积分操作都是绝对合法的。例如对上式两边进行逐项积分：
\[
\int_0^x\frac{dt}{1-t}=\sum_{n=0}^{\infty}\int_0^x t^n\,dt
\quad\Longrightarrow\quad
-\ln(1-x)=\sum_{n=1}^{\infty}\frac{x^n}{n},\qquad\abs x<1.
\]

注意：微积分操作导出的结论\textbf{默认只在区间内部 \(|x| < 1\) 保真}。内部的操作不会自动延伸到端点。至于 \(x=-1\) 处是否能成立 \(\ln 2=1-\frac12+\frac13-\cdots\)，需要靠专门的端点连续性定理（如 Abel 定理）或余项估计来补齐证明。

\subsection{为什么内部可以任意逐项做微积分？}

\begin{theorem}
设 \(f(x)=\sum_{n\ge0}c_n(x-a)^n\) 的收敛半径为 \(R>0\)。则 \(f(x)\) 在开区间 \((a-R,a+R)\) 内无穷次可微，且可以任意次逐项求导与逐项积分，得到的新幂级数\textbf{收敛半径依然是 \(R\)}。
\end{theorem}

这个结论的威力比表面看起来要可怕得多。在高等数学中，``无穷项求和'与``求导/积分'交换顺序是极度危险的操作，极易导致错误。

幂级数之所以能在收敛区间内部畅通无阻，靠的是内部存在着极其充沛的\textbf{``安全缓冲垫'}：在任何严格包含于内部的闭区间 \([a-r, a+r]\)（其中 \(r < R\)）上，尾部项都可以被一个收敛的几何级数\textbf{统一且一致地压制住}（即一致收敛）。

不过请注意：虽然微积分不改变收敛半径 \(R\)，但\textbf{可能会改变端点的生死状态}：
\begin{itemize}
    \item \textbf{求导}：系数乘以 \(n\)，尾巴变重，倾向于把原本收敛的端点``拉下水'（使其发散）；
    \item \textbf{积分}：系数除以 \(n\)，尾巴变轻，倾向于把原本发散的端点``救回来'（使其收敛）。
\end{itemize}
例如 \(\sum x^n\) 在 \(\pm 1\) 均发散，积分一次变成 \(\sum \frac{x^n}{n}\)，在 \(x=-1\) 处就奇迹般地活了过来。因此，\textbf{每次做完微积分，端点行为都必须重新查验}。

\subsection{工程陷阱：数学上收敛 \(\neq\) 计算上可用}

最后，必须揭示一个在实际工程计算与算法设计中经常踩坑的落差：\textbf{数学上的收敛保证，并不等于工程上的可计算性。}

理论上的收敛半径只保证``只要取足够多的项，误差就可以任意小'，但它\textbf{从来没保证过``足够多'到底是多少项}。

以 \(-\ln(1-x) = \sum_{n=1}^{\infty} \frac{x^n}{n}\) 为例：
在 \(x = 0.99\) 处，它确实位于收敛区间 \(|x| < 1\) 内部，数学上绝对收敛。然而，它的项按 \(\frac{0.99^n}{n}\) 缓慢衰减。如果要将截断误差控制在 \(10^{-6}\) 以内，你需要累加上千项！
更糟的是，上千个浮点数在计算机里连续相加，累积的舍入误差可能会直接掩盖掉后几项的真值。

\textbf{在靠近收敛半径边界的地方，级数在数学上完全有效，但在工程计算上几乎废柴。}

这就是为什么工业级数学库（如计算 \(\ln, \sin, e^x\) 的底层算法）从来不会直接把远处的 \(x\) 代入展开式。工程上的标准做法是：先通过函数的代数恒等式（如周期性、倍角公式或区间归约），将 \(x\) 强行\textbf{``拉回'到展开中心 \(a\) 的极小邻域内}，再用区区 \(3 \sim 5\) 项极快速地计算出极其精确的结果。

展开中心选在哪里、截断保留几项、由此带来多大的浮点与截断误差，这三者在工程落地时必须联合计算。下一节，我们将正式把系数从``给定'换成``由函数的各阶导数生成'，并彻底解决那个最关键的问题：\textbf{截断之后，到底差了多少？}
